\documentclass[10pt]{beamer}

% --- Configurazione Tema ---
\usetheme{Madrid}
\usecolortheme{default}

% --- Pacchetti ---
\usepackage[utf8]{inputenc}
\usepackage[T1]{fontenc}
\usepackage[italian]{babel}
\usepackage{tikz}
\usetikzlibrary{arrows.meta, positioning, shapes.geometric, decorations.pathreplacing, calc, shadows, fit}
\usepackage{listings}
\usepackage{xcolor}
\usepackage{booktabs}

% --- Colori Personalizzati ---
\definecolor{primary}{RGB}{38, 66, 139}
\definecolor{codebg}{RGB}{245, 245, 245}
\definecolor{keywordcolor}{RGB}{0, 0, 255}
\definecolor{stringcolor}{RGB}{163, 21, 21}
\definecolor{commentcolor}{RGB}{0, 128, 0}

\setbeamercolor{structure}{fg=primary}
\setbeamercolor{title}{fg=white, bg=primary}

% --- Configurazione Codice ---
\lstdefinelanguage{json}{
    basicstyle=\ttfamily\scriptsize,
    showstringspaces=false,
    breaklines=true,
    frame=lines,
    backgroundcolor=\color{codebg},
    string=[s]{"}{"},
    stringstyle=\color{stringcolor},
    keywordstyle=\color{keywordcolor},
}

\lstdefinelanguage{javascript}{
    keywords={async, await, break, case, catch, class, const, continue, debugger, default, delete, do, else, export, extends, false, finally, for, function, if, import, in, instanceof, new, null, return, super, switch, this, throw, true, try, typeof, var, void, while, with, yield, let},
    morekeywords={console, log, Error, fetch, JSON, parse, stringify},
    sensitive=true,
    morecomment=[l]{//},
    morecomment=[s]{/*}{*/},
    morestring=[b]',
    morestring=[b]",
    morestring=[b]`,
    keywordstyle=\color{primary},
    commentstyle=\color{commentcolor},
    stringstyle=\color{stringcolor},
    basicstyle=\ttfamily\scriptsize
}

\lstset{
    backgroundcolor=\color{codebg},
    basicstyle=\ttfamily\scriptsize,
    breaklines=true,
    frame=single,
    keywordstyle=\color{keywordcolor},
    commentstyle=\color{commentcolor},
    stringstyle=\color{stringcolor}
}

% --- Informazioni ---
\title[Backend e API]{Programmazione e Tecnologie Web}
\subtitle{Lezione 3: Il Backend e le API RESTful}
\author{Giuseppe Capasso}
\institute[ITS]{Istituto Tecnico Superiore (ITS)}
\date{A.A. 2025/2026}

\begin{document}

% --- Slide Titolo ---
\begin{frame}
    \titlepage
\end{frame}

% --- Obiettivi ---
\begin{frame}{Obiettivi della Lezione}
    \begin{itemize}
        \item \textbf{Architetture Moderne}: Il paradigma API-First e Microservizi.
        \item \textbf{RESTful Design}: Semantica HTTP, JSON e Stato.
        \item \textbf{Cybersecurity}: AAA, JWT e OAuth 2.0.
        \item \textbf{Design Patterns}: Versioning, Filtering e Paginazione.
        \item \textbf{Laboratorio}: Consumo API con Python e Node.js.
    \end{itemize}
\end{frame}

% ==========================================
% CAPITOLO 1: ARCHITETTURE
% ==========================================
\section{Architetture API-First}

\begin{frame}{Il Paradigma API-First}
    \begin{itemize}
        \item Il backend non genera HTML (SSR), ma espone \textbf{dati grezzi}.
        \item L'API \`e il "contratto" tra client e server.
    \end{itemize}
    \centering
    \resizebox{0.7\textwidth}{!}{
    \begin{tikzpicture}[
        node distance=1.5cm,
        client/.style={rectangle, draw, rounded corners, fill=blue!5, minimum width=2.5cm, minimum height=1cm, align=center, font=\small},
        server/.style={rectangle, draw, rounded corners, fill=primary!20, minimum width=3.5cm, minimum height=1.5cm, align=center, font=\bfseries},
        database/.style={cylinder, draw, rotate=90, fill=gray!10, minimum width=1.2cm, minimum height=1.5cm, align=center, font=\small}
    ]
        \node (web) [client] {SPA Client\\(React/Angular)};
        \node (mobile) [client, below=of web] {Mobile Client\\(iOS/Android)};
        \node (iot) [client, below=of mobile] {IoT/CLI\\(Python/Node)};
        
        \node (api) [server, right=3cm of mobile] {API Gateway /\\Backend Service};
        \node (db) [database, right=2.5cm of api] {\rotatebox{-90}{Database}};

        \draw [latex-latex, line width=1pt, primary] (web.east) -- node[above, black, scale=0.7, sloped] {HTTP/JSON} (api.west);
        \draw [latex-latex, line width=1pt, primary] (mobile.east) -- node[above, black, scale=0.7] {HTTP/JSON} (api.west);
        \draw [latex-latex, line width=1pt, primary] (iot.east) -- node[below, black, scale=0.7, sloped] {HTTP/JSON} (api.west);
        \draw [dashed, line width=0.8pt] (api.east) -- (db.north);
    \end{tikzpicture}
    }
\end{frame}

\begin{frame}{Monoliti vs Microservizi}
    \centering
    \resizebox{0.8\textwidth}{!}{
    \begin{tikzpicture}[
        node distance=0.8cm and 1.5cm,
        module/.style={rectangle, draw, fill=gray!5, minimum width=2.2cm, minimum height=0.8cm, font=\scriptsize, align=center},
        service/.style={rectangle, draw, rounded corners, fill=primary!10, minimum width=2.5cm, minimum height=1cm, font=\small, align=center},
        mono_container/.style={rectangle, draw, thick, dashed, inner sep=0.4cm, fill=gray!2, label=above:\textbf{Monolith}},
        gateway/.style={trapezium, draw, trapezium left angle=120, trapezium right angle=120, fill=primary!20, minimum width=4cm, font=\small\bfseries, align=center}
    ]
        % Monolith
        \node (m1) [module] {Auth Module};
        \node (m2) [module, below=of m1] {User Module};
        \node (m3) [module, below=of m2] {Order Module};
        \node (m4) [module, below=of m3] {Notify Module};
        \node [mono_container, fit=(m1) (m4)] {};

        % Microservices
        \node (gw) [gateway, right=5cm of m1, yshift=-1.5cm] {API Gateway};
        \node (s1) [service, right=2cm of gw, yshift=1.5cm] {Auth Service};
        \node (s2) [service, below=0.5cm of s1] {User Service};
        \node (s3) [service, below=0.5cm of s2] {Order Service};
        \node (s4) [service, below=0.5cm of s3] {Notify Service};

        \draw [latex-latex, primary] (gw.east) -- (s1.west);
        \draw [latex-latex, primary] (gw.east) -- (s2.west);
        \draw [latex-latex, primary] (gw.east) -- (s3.west);
        \draw [latex-latex, primary] (gw.east) -- (s4.west);
    \end{tikzpicture}
    }
    \vspace{0.3cm}
    \begin{columns}
        \begin{column}{0.5\textwidth}
            \textbf{Monolito}: Semplice, ma difficile da scalare e mantenere.
        \end{column}
        \begin{column}{0.5\textwidth}
            \textbf{Microservizi}: Complesso, ma scalabile e resiliente.
        \end{column}
    \end{columns}
\end{frame}

\begin{frame}{Tassonomia delle API}
    \centering
    \resizebox{0.7\textwidth}{!}{
    \begin{tikzpicture}[
        layer/.style={rectangle, draw, thick, minimum width=8cm, minimum height=1.2cm, align=center, font=\small\bfseries},
        arrow/.style={-latex, line width=1pt, primary}
    ]
        \node (web) [layer, fill=primary!20] {Web APIs (REST, GraphQL, RPC)\\ \normalfont\scriptsize Comunicazione via Rete (HTTP)};
        \node (browser) [layer, fill=primary!10, below=0.5cm of web] {Browser APIs (DOM, Fetch, Canvas)\\ \normalfont\scriptsize Runtime del Client};
        \node (lib) [layer, fill=gray!10, below=0.5cm of browser] {Library APIs (Standard Library, SDKs)\\ \normalfont\scriptsize Spazio di Memoria Locale};

        \draw [arrow] (browser.north) -- (web.south);
        \draw [arrow] (lib.north) -- (browser.south);
    \end{tikzpicture}
    }
\end{frame}

% ==========================================
% CAPITOLO 2: REST & JSON
% ==========================================
\section{REST e Scambio Dati}

\begin{frame}[fragile]{Serializzazione e JSON}
    \centering
    \begin{tikzpicture}[
        box/.style={rectangle, draw, rounded corners, minimum width=3cm, minimum height=1.2cm, align=center, font=\small},
        arrow/.style={-latex, line width=1pt, primary}
    ]
        \node (mem) [box, fill=gray!10] {Oggetto in Memoria\\(Dict / Object)};
        \node (wire) [box, fill=primary!10, right=4cm of mem] {Stringa JSON\\(Payload Rete)};

        \draw [arrow] ([yshift=0.2cm]mem.east) -- node[above, scale=0.7] {Serializzazione} ([yshift=0.2cm]wire.west);
        \draw [arrow] ([yshift=-0.2cm]wire.west) -- node[below, scale=0.7] {Deserializzazione} ([yshift=-0.2cm]mem.east);
    \end{tikzpicture}
    
    \vspace{0.3cm}
    \begin{lstlisting}[language=json]
{
  "id": 42,
  "title": "Introduzione alle API",
  "author": {"name": "Giuseppe Capasso"},
  "tags": ["web", "api", "rest"]
}
    \end{lstlisting}
\end{frame}

\begin{frame}{Semantica HTTP in REST}
    \centering
    \begin{table}
    \begin{tabular}{lccl}
    \toprule
    \textbf{Metodo} & \textbf{Sicuro} & \textbf{Idempotente} & \textbf{Azione REST} \\
    \midrule
    \texttt{GET}    & S\`i & S\`i & Legge una risorsa \\
    \texttt{POST}   & No & No & Crea una risorsa \\
    \texttt{PUT}    & No & S\`i & Sostituisce risorsa \\
    \texttt{PATCH}  & No & No & Modifica parziale \\
    \texttt{DELETE} & No & S\`i & Elimina risorsa \\
    \bottomrule
    \end{tabular}
    \end{table}
    \vspace{0.3cm}
    \begin{itemize}
        \item \textbf{Sicurezza}: Non modifica lo stato sul server.
        \item \textbf{Idempotenza}: Pi\`u chiamate identiche = stesso effetto di una.
    \end{itemize}
\end{frame}

\begin{frame}{HTTP Status Codes}
    \begin{itemize}
        \item \textbf{2xx (Successo)}:
        \begin{itemize}
            \item \texttt{200 OK}, \texttt{201 Created}, \texttt{204 No Content}.
        \end{itemize}
        \item \textbf{4xx (Errore Client)}:
        \begin{itemize}
            \item \texttt{400 Bad Request}: Dati inviati errati.
            \item \texttt{401 Unauthorized}: Autenticazione mancante.
            \item \texttt{403 Forbidden}: Mancanza di permessi.
            \item \texttt{404 Not Found}: Risorsa inesistente.
        \end{itemize}
        \item \textbf{5xx (Errore Server)}:
        \begin{itemize}
            \item \texttt{500 Internal Server Error}.
        \end{itemize}
    \end{itemize}
\end{frame}

% ==========================================
% CAPITOLO 3: SICUREZZA
% ==========================================
\section{Sicurezza delle API}

\begin{frame}{Il Modello AAA}
    \begin{enumerate}
        \item \textbf{Authentication} (Chi sei?): Verifica identit\`a.
        \item \textbf{Authorization} (Cosa puoi fare?): Controllo permessi.
        \item \textbf{Accounting} (Cosa hai fatto?): Logging e Audit.
    \end{enumerate}
    \vspace{0.5cm}
    \textbf{Fattori di Autenticazione}:
    \begin{itemize}
        \item \textbf{Something you know}: Password, PIN.
        \item \textbf{Something you have}: Token hardware, Smartphone (OTP).
        \item \textbf{Something you are}: Biometria (Impronta, Viso).
    \end{itemize}
\end{frame}

\begin{frame}{JWT: JSON Web Tokens}
    \centering
    \resizebox{0.8\textwidth}{!}{
    \begin{tikzpicture}[
        node distance=0cm,
        part/.style={rectangle, draw, minimum width=3cm, minimum height=1cm, align=center, font=\ttfamily\small, thick},
        header_style/.style={fill=red!15},
        payload_style/.style={fill=blue!15},
        sig_style/.style={fill=green!15}
    ]
        \node (header) [part, header_style] {Header};
        \node (p1) [right=0.2cm of header, font=\huge] {.};
        \node (payload) [part, payload_style, right=0.2cm of p1] {Payload};
        \node (p2) [right=0.2cm of payload, font=\huge] {.};
        \node (sig) [part, sig_style, right=0.2cm of p2] {Signature};

        \node [below=0.2cm of header, font=\scriptsize, text=red!70!black] {Algoritmo \& Tipo};
        \node [below=0.2cm of payload, font=\scriptsize, text=blue!70!black] {Claims (Dati Utente)};
        \node [below=0.2cm of sig, font=\scriptsize, text=green!70!black] {Firma Integrit\`a};
    \end{tikzpicture}
    }
    \vspace{0.5cm}
    \begin{itemize}
        \item \textbf{Stateless}: Il server non deve salvare sessioni nel DB.
        \item \textbf{Self-contained}: Contiene tutte le info necessarie.
    \end{itemize}
\end{frame}

\begin{frame}{OAuth 2.0: Authorization Code Flow}
    \centering
    \resizebox{0.7\textwidth}{!}{
    \begin{tikzpicture}[
        node distance=1.5cm, font=\sffamily\small,
        line/.style={draw, -latex, thick, primary},
        lifeline/.style={draw, dashed, thick, gray!50},
        actor/.style={rectangle, draw, rounded corners, fill=primary!10, minimum width=2cm, minimum height=0.7cm, align=center, font=\bfseries\scriptsize}
    ]
        \node (user) [actor] {User};
        \node (client) [actor, right=2cm of user] {Client App};
        \node (auth) [actor, right=2cm of client] {Auth Server};
        \node (api) [actor, right=2cm of auth] {API Server};

        \draw [lifeline] (user.south) -- ++(0,-5.5);
        \draw [lifeline] (client.south) -- ++(0,-5.5);
        \draw [lifeline] (auth.south) -- ++(0,-5.5);
        \draw [lifeline] (api.south) -- ++(0,-5.5);

        \draw [line] ($(user.south)-(0,0.5)$) -- node[above, scale=0.6] {1. Login Request} ($(client.south)-(0,0.5)$);
        \draw [line] ($(client.south)-(0,1.2)$) -- node[above, scale=0.6, sloped] {2. Redirect to Auth} ($(auth.south)-(0,1.2)$);
        \draw [line] ($(auth.south)-(0,1.9)$) -- node[above, scale=0.6] {3. Auth Code} ($(client.south)-(0,1.9)$);
        \draw [line] ($(client.south)-(0,2.6)$) -- node[above, scale=0.6] {4. Exchange Code} ($(auth.south)-(0,2.6)$);
        \draw [line] ($(auth.south)-(0,3.3)$) -- node[above, scale=0.6] {5. Access Token} ($(client.south)-(0,3.3)$);
        \draw [line] ($(client.south)-(0,4.0)$) -- node[above, scale=0.6] {6. Call API + Token} ($(api.south)-(0,4.0)$);
    \end{tikzpicture}
    }
\end{frame}

% ==========================================
% CAPITOLO 4: DESIGN PATTERNS
% ==========================================
\section{Design Patterns Avanzati}

\begin{frame}{Paginazione: Offset vs Cursor}
    \centering
    \resizebox{0.8\textwidth}{!}{
    \begin{tikzpicture}[
        record/.style={rectangle, draw, minimum width=0.8cm, minimum height=0.6cm, fill=gray!10},
        active_record/.style={rectangle, draw, minimum width=0.8cm, minimum height=0.6cm, fill=primary!20, thick},
        arrow/.style={-latex, line width=1pt, primary}
    ]
        % Offset
        \node (o1) [record] {1}; \node (o2) [record, right=0.1cm of o1] {2}; \node (o3) [record, right=0.1cm of o2] {3};
        \node (o4) [active_record, right=0.1cm of o3] {4}; \node (o5) [active_record, right=0.1cm of o4] {5};
        \draw [decorate, decoration={brace, amplitude=5pt}] (o1.north west) -- (o3.north east) node [black, midway, yshift=0.4cm, font=\scriptsize] {Offset (skip 3)};

        % Cursor
        \begin{scope}[yshift=-1.5cm]
            \node (c3) [record, label=below:{\scriptsize Cursor}] {3};
            \node (c4) [active_record, right=0.1cm of c3] {4}; \node (c5) [active_record, right=0.1cm of c4] {5};
            \draw [arrow] (c3.north) -- ++(0,0.3) -- ++(0.9,0) -- (c4.north);
        \end{scope}
    \end{tikzpicture}
    }
    \vspace{0.3cm}
    \begin{itemize}
        \item \textbf{Offset}: Facile, ma lento su milioni di righe.
        \item \textbf{Cursor}: Performante e stabile su dati real-time.
    \end{itemize}
\end{frame}

% ==========================================
% CAPITOLO 5: LABORATORIO
% ==========================================
\section{Laboratorio Pratico}

\begin{frame}[fragile]{Python: requests}
    \begin{lstlisting}[language=python]
import requests
url = "https://jsonplaceholder.typicode.com/posts/1"

try:
    response = requests.get(url, timeout=5)
    response.raise_for_status() 
    
    data = response.json()
    print(f"Titolo: {data['title']}")

except Exception as e:
    print(f"Errore: {e}")
    \end{lstlisting}
    \textbf{Esecuzione}: \texttt{python client.py}
\end{frame}

\begin{frame}[fragile]{Node.js: fetch (v18+)}
    \begin{lstlisting}[language=javascript]
async function getPost() {
    const url = "https://jsonplaceholder.typicode.com/posts/1";
    try {
        const response = await fetch(url);
        if (!response.ok) throw new Error("Errore HTTP");

        const data = await response.json();
        console.log("Titolo:", data.title);
    } catch (error) {
        console.error("Errore:", error.message);
    }
}
getPost();
    \end{lstlisting}
    \textbf{Esecuzione}: \texttt{node api\_test.js}
\end{frame}

\begin{frame}{Conclusioni}
    \begin{itemize}
        \item Le API sono il \textbf{tessuto connettivo} del web moderno.
        \item Progettare bene significa pensare a \textbf{sicurezza} e \textbf{scalabilit\`a} sin dal primo giorno.
        \item \textbf{Prossimi Passi}: Sperimentare con OpenAPI (Swagger) per documentare i vostri servizi.
    \end{itemize}
    \centering
    \vspace{0.8cm}
    \Large \textbf{\textcolor{primary}{Domande?}}
\end{frame}

\end{document}
